\label{sec:north-carolina}

\subsection{Background: \textit{Harper v.\ Hall}}

\textit{Harper v.\ Hall}~\citep{harper2022} is the most significant state
court partisan gerrymandering decision since \textit{Rucho v.\ Common
Cause}~\citep{rucho2019} closed the federal courthouse door.  The North
Carolina Supreme Court held, 4-3, that the North Carolina General Assembly's
enacted congressional and state legislative maps violated the state
constitution's free elections clause and equal protection guarantee.  The
court ordered the legislature to redraw the maps and retained jurisdiction to
review the remedial plan.

\textbf{Note on subsequent procedural history.}
Following the 2022 election, the North Carolina Supreme Court's reconstituted
5-2 Republican majority reversed its earlier \textit{Harper v.\ Hall} holding
in April 2023, reinstating the legislature-drawn map.
This case study describes the remedial-map scenario as an illustration of
how a court \emph{could} use algorithmic redistricting, not as the actual
outcome of \textit{Harper}.
The \textit{Harper} fact pattern nevertheless established the template for
state-court-ordered redistricting remedies in the partisan gerrymandering
context, and the algorithmic approach demonstrated here would apply to any
future court order in North Carolina or any other state with a
partisan-gerrymandering clause in its state constitution.

\subsection{The Special Master Role}

When a court orders a redistricting remedy, it typically has three options:
(a) order the legislature to redraw the map under specified criteria;
(b) draw the map itself; or (c) appoint a special master to draw the map.
Option (a) risks delay and noncompliance.  Option (b) is constitutionally
problematic as courts are not representative bodies.  Option (c) has
become the standard remedy in federal redistricting cases and is available
to state courts as well.

The algorithmic redistricting system is particularly well-suited to the
special master context.  A special master appointed by the court could be
instructed to:

\begin{enumerate}
  \item Run \texttt{redist state --state NC --year 2020 --version remedy}
    with population-balance and compactness objectives.
  \item Apply VRASection weights to protect the Black-majority district
    in the northeastern part of the state (the historical Gingles area)
    and any other minority-opportunity district required by VRA analysis.
  \item Submit the algorithm output to the court as the proposed remedial
    plan, with a criteria-compliance report generated by
    \texttt{redist analyze --state NC --year 2020 --version remedy
    --types compactness demographic political}.
  \item Defend the plan in court proceedings, addressing any constitutional
    challenge to the output.
\end{enumerate}

\subsection{The Algorithm Output as the Presumptive Remedy}

The court-order pathway suggests a strong procedural innovation: the algorithm
output should be presumed constitutional absent a specific showing of
constitutional defect.  This presumption reflects two structural features
of algorithmic redistricting.

First, the algorithm cannot gerrymander because it cannot see partisan data.
A constitutional challenge on partisan-gerrymandering grounds would need to
show that the \emph{output} constitutes a constitutional violation regardless
of the process---a much harder showing than demonstrating partisan intent.

Second, the algorithm satisfies the neutral criteria (population equality,
compactness) that every court has endorsed as constitutional.  A challenger
must therefore show not merely that a different map would be more favorable
but that the algorithm's output violates a specific constitutional requirement.

This structure shifts the burden of proof in a legally significant way.
Rather than the state defending a legislatively drawn map against allegations
of partisan manipulation, the special master defends an algorithmically drawn
map against allegations of specific constitutional defects.  The latter is a
substantially easier task.

\subsection{Scope of Judicial Review}

The court-order pathway does not eliminate judicial review; it changes its
focus.  Judicial review of an algorithmic remedial map would appropriately
address:
\begin{itemize}
  \item Whether the algorithm's population-balance constraint is adequate
    (the equal-population requirement).
  \item Whether the VRASection weighting adequately protects minority-opportunity
    districts (VRA compliance).
  \item Whether the algorithm's compactness objective is consistent with
    state constitutional criteria (state law compliance).
  \item Whether any community-of-interest exceptions identified by the special
    master are supported by factual findings.
\end{itemize}

Questions of partisan outcome---which party benefits, what the expected seat
count is---would be legally irrelevant under \textit{Rucho} in the federal
context and, after \textit{Harper}, would need to be evaluated against the
specific constitutional standard adopted by the state court.  Crucially,
because the algorithm does not access partisan data, no partisan intent can be
attributed to the mapmaker; partisan-outcome challenges must rest on effects
rather than intent.
